\scanpage{76}\mspage{K}
\uline{Théorème} Soit $K$ une extension de type fini du
corps $k$, \deleted{[ayant au moins un modèle
propre \inserted{\unsure{int.}} et universellement normal sur $k$]} \deleted{Soit} Il
\deleted{$(X_i)_{i\in\unsure{I|+}}$ le système ind} existe un entier $n(K/k)$
tel que pour \deleted{tout} modèle \inserted{$X$ de $K$} propre sur $k$, et
normal sur $\unsure{k|K}$, \deleted{\unsure{pour}} on ait \deleted{$\dim\underline{\mathrm{Pic}}_X$}
\[
\dim\underline{\mathrm{Pic}}_{X/k}\leq n(K/k)
\]

Démonstration. \deleted{On est ramené au cas} \inserted{S'il n'existe pas de}
modèle $X$ du type indiqué, c'est trivial.
Sinon, $K$ est séparable sur $k$, et
$\overline K=K\otimes_k\bar k$ est une extension de type
fini du corps alg. clos $\unsure{\bar K|\bar k}$, et on
est ramené au cas où $\bar k$ est alg. clos.
On a alors
\[
\begin{aligned}
\dim\underline{\mathrm{Pic}}_{X/k}
&=\tfrac12\operatorname{rang}_{\mathbb Z_{\ell}}
  \bigl(\text{partie $l$-primaire}\\
&\qquad \text{de $\pi_1(X/k)$ rendu}\\
&\qquad \text{abélien}\bigr)
\end{aligned}
\]
et à ce titre, en vertu du théorème
de majoration uniforme \unsure{des|du} $\pi_1(X/k)$
% archive sheet number: 75
